\chapter{1965:Robert B. Woodward}

\label{chap:chemistry1965}

The Nobel Prize in Chemistry for the year 1965 was awarded to Robert B. Woodward.

\textbf{Motivation:}

for his outstanding achievements in the art of organic synthesis

\section{Robert B. Woodward}

\subsection*{Early Life and Education}

Robert B. Woodward was born on 1917-04-10 in Boston, Massachusetts, USA.

\subsection*{Career and Major Research}

Woodward studied at the Massachusetts Institute of Technology, where he received his doctorate in 1937, and joined Harvard University, where he was professor of chemistry from 1950. He completed the total synthesis of quinine in 1944 and later led teams that synthesized cholesterol, cortisone, strychnine, chlorophyll, and vitamin B12. Together with Roald Hoffmann he formulated the Woodward-Hoffmann rules on the conservation of orbital symmetry.

\subsection*{Nobel Prize-winning Work}

The work for which Robert B. Woodward was awarded the Nobel Prize in Chemistry in 1965 was described by the Nobel Committee as: "for his outstanding achievements in the art of organic synthesis".

Woodward's syntheses of complex natural products, beginning with quinine in 1944, established new standards for planning and executing multistep syntheses. His approach combined rigorous structural reasoning with carefully designed reaction sequences and physical methods of analysis.

\subsection*{Significance and Impact}

The total syntheses of quinine, cholesterol, cortisone, strychnine, chlorophyll, and vitamin B12 demonstrated that complex natural molecules could be built from simple starting materials. These achievements made him one of the most influential organic chemists of the twentieth century and shaped modern synthetic methodology.

\subsection*{Later Career and Honors}

Woodward remained at Harvard for his entire academic career and also served as a director of research at the Woodward Research Institute in Basel. He died on 1979-07-08 in Cambridge, Massachusetts.

\subsection*{Legacy}

The Woodward-Hoffmann rules and Woodward's total syntheses remain cornerstones of organic chemistry education and research. The many students trained in his laboratory carried his approach into the wider field.

\subsection*{Modern Developments}

Woodward's total syntheses, including cholesterol, chlorophyll, and vitamin B12, set the standard for organic synthesis and founded modern synthetic chemistry. His ideas, formalized in the Woodward–Hoffmann rules, guide the design of countless reactions used in the pharmaceutical and agrochemical industries. Woodward's approach to synthesis remains the blueprint for complex molecule construction today.

\subsection*{Selected Publications}

"The Total Synthesis of Quinine" (Journal of the American Chemical Society, 1944)

\clearpage

\section*{Chapter Summary}

The 1965 prize recognized Robert B. Woodward for his total syntheses of complex natural products, from quinine in 1944 to vitamin B12. His work established the standards of modern synthetic organic chemistry.
